\section{Draft Notes}

Marco had this idea that casts may be avoided by using dynamic dispatch returning this.
Can not work, since Meth local class is private, there is no way an external interface can mention it to provide the cast method.
Of course, the method local is private to the body (this means also the method signature can not use it)
This means no self:Foo method, nor self:T method of a generic class eventually instantiating T=Foo.

This means the core language DO NEED if(e instanceof T x){} else {}

//nested funnelling:
  new A<X>(){ B<X> f()-> new B<X>(){} }
  new A<X>(){ <K>B<X,K> f()-> new B<X,K>(){} }

using 'extends' for declarations in top level classes and generic methods
using 'funnel' for named inner classes
  new A<X funnel Object>(){ B<X> f()-> new B<X funnel Object>(){} }
  new A<X funnel Object>(){ <K extends Object> B<X,K> f()-> new B<X funnel Object,K funnel Object>(){} }
  or, if the bounds are not repeated
  new A<X funnel>(){ <K extends Object> B<X,K> f()-> new B<X funnel,K funnel>(){} }
Very funky syntax:
we get new B<X funnel Object,K funnel Object>(){}
but, during reduction, this can become 
we get new B<Person,K funnel Object>(){}
where we have a mixture of concrete types and generic funnelling

Global class definition uses class C<X extends Bar>{..}
named method local class uses new C<X funnel>(){..}
normal new: new C<Ts>()

\begin{verbatim}
e ::= x | e.m(es) | e.f | new C<Ts>(es) | new C<Fs>(){Ms} | if (e0 instanceof T x){ e1 }else{ e2 }
T ::= C<Ts> | X
F ::= C<CCs> | X funnel// Important, can not be T | X funnel, otherwise we could have an X not 'funnelled'
CC::= C<CCs> //types with zero free type variables
B::=  X extends T // Need T, not F, to be able to do things like (Java) class Bar<X,Y extends X> {}
Dec::= class C<Bs>{ Ms Ss} //assume standard constructor
S ::= T x; //field/slot
M ::= Sig; | Sig{ return e; }
Sig::= <Bs> T m(x1:T1,..xn:Tn)
G ::= x1:T1..xn:Tn
\end{verbatim}
//we do not need Object to be able to write 'extends Object'
//this is exactly as the distinction of checked and unchecked exceptions: formally it evaporates in syntacitc sugar where all is checked
//and all methods 'throws Error, RTE'
// Can assume the user provides their own Object type
Note: an use of X is just X, a declaration of X is X extends T; a funnelling of X is X funnel
Program is just many Dec plus an e typable in empty
MARCO TODO: 
  typechecking
    Decs;Bs;G|-e:T
    Decs|-Dec:ok
    Decs;T|-M:ok
    e --> e'
Note, there are multiple semantics/ts:
  e instanceof CC //always clear
  e instanceof T with T not of form CC would not typecheck

  e --> e' for 
  Java with warn->err: just have CC in instanceof
  Java with magic reifined cast: just have T in instanceof (since reduction will replace(thus remember) generics
  Java with warn->off:
        AT RUN TIME instanceof only checks base class, ignores generic parameters
        AT COMPLILE TIME, instanceof can be err if the full type (including generics) is not potentially correlated)?? unobvious consider interfaces
        In prose, should explain that we won't flag cast of some type Foo to some type Bar, even though Java sometimes flags unrelated casts
\mjnote{these three semantics are replaced by the single one of Java 16+: erased run-time check, plus the checked-cast premise of rule \textsc{if-instanceof} (Section~\ref{sec:cast}). Unrelated types are rejected by it, as in Java.}

<X extends List<Y funnel>>

class Foo<X> {
    class Getter<X funnel> {
        <Y extends List<X>> X getFirst(Y y) { y.getOrElse(0,null) }
    }
}

<Y> .. return foo<Y,MyX>();
<Y,X extends List<Y>> foo(){ return Foo<X funnel>(){... X }; }
